\section{Introduction}
\label{sec:intro}

``The intervention did not change the behavior'' is one of the most common sentences in alignment research, and one of the hardest to interpret.
A steering vector that leaves outputs unchanged, a system prompt that fails to reduce sycophancy, an ablation with no measurable effect: each of these is a null result, and each is read in two opposite ways.
One reading says the behavior is robust to this kind of intervention.
The other says the experiment failed to detect something that was there.
The two readings lead to opposite safety conclusions, so it matters which one the evidence supports.

A null result is the joint product of a model and an experimental design.
The same intervention on the same model can be scored by whether the sampled answer changes, whether the most likely answer changes, or whether the answer logits move; it can be tested on 20 items or 200.
If these choices largely fix whether a study returns a null, then a null says less about the model than it appears to.
Other fields have met this problem before.
Multiverse and specification-curve analyses in psychology run every defensible analysis of one dataset and report how the conclusion varies~\cite{steegen2016increasing,simonsohn2020specification}.
Work on evaluation statistics for language models addresses sample size and variance~\cite{miller2024adding,card2020little}, and work on steering vectors shows that their effects vary widely across inputs~\cite{tan2024analyzing}.
However, we know of no study that takes one fixed set of alignment interventions through a full grid of design choices and maps where the nulls fall.

We build such a map for one behavior, sycophancy (see \figref{fig:overview}).
We run eleven interventions at four levels of granularity (prompt, all-layer activation, single-layer activation, and attention head), including a placebo at each level, on four \qwen models from 0.5B to 7B parameters~\cite{qwen2024qwen25}, using the same 240 evaluation items from the model-written sycophancy evaluations~\cite{perez2022discovering}.
We then treat each of the 44 model--intervention pairs as a population from which we draw simulated studies: 1000 resamples for each of four metrics and four sample sizes, 704 design cells in total.
The fraction of studies in a cell with $p \ge .05$ is that cell's null rate.
Because every cell uses the same underlying items, we can also estimate a reference effect for each intervention on the full pool and ask what the nulls are made of.

Our hypothesis has two parts.
The first is that nulls cluster along design features: metric, intervention granularity, and model scale.
The second is that this clustering reflects a methodological blind spot, so that most nulls occur where a real effect exists.
We fixed one asymmetry before collecting data.
Metric and sample size are measurement-side features: the underlying effect is identical across their levels, so any difference in null rate is methodological by construction.
Granularity and scale change the intervention or the model, so clustering there could reflect different true effects.

{\bf What do we find?}
Nulls cluster strongly, but on the measurement side.
The null rate is 0.81 for a sampled answer, 0.80 for the greedy answer, 0.52 for answer probability, and 0.25 for the logit margin (within-unit permutation $p = 0.0002$).
Moving from an answer-level metric to the logit margin changes the null rate by 0.56, more than moving from 20 to 200 items (0.34).
Granularity does not organize the nulls (permutation $p = 0.88$), and model scale has no monotone effect.
The second part of the hypothesis is not supported as we pre-specified it.
Only 14\% (95\% CI 5--27\%) of nulls from non-placebo interventions fall on interventions that demonstrably change the model's answers.
Instead, 80\% (62--93\%) fall on interventions that shift the answer logits without reliably flipping answers, and 6\% fall on the single intervention with no detectable effect.
The typical null in our study is neither a missed behavioral effect nor an inert intervention: the intervention moves the model internally by an amount too small to change what it says.

In summary, our contributions are:
\begin{itemize}[leftmargin=*,itemsep=0pt,topsep=0pt]
    \item We propose a controlled multiverse design for alignment interventions that holds model, intervention, and items fixed while varying metric and sample size, and pairs each design cell with a full-pool reference effect.
    \item We map null rates over 704 design cells built from real outputs of four models under eleven sycophancy interventions, and show that metric choice explains 32\% of the variance in null rate while granularity explains 2\%.
    \item We classify every null by the reference effect of its intervention and show that 80\% of non-placebo nulls are logit-level shifts that leave answers unchanged, which neither the ``robustness'' nor the ``blind spot'' reading describes well.
    \item We document three results that complicate standard practice: placebo interventions are not inert, attribution-selected heads can do less than random heads, and logit-level and answer-level metrics can disagree in sign.
\end{itemize}

The evidence covers one behavior, one model family, and a forced-choice format; we discuss these limits in \secref{sec:discussion}.
\Secref{sec:related} reviews related work, \secref{sec:method} describes the design, \secref{sec:results} reports results, and \secref{sec:conclusion} concludes.
